% ECONOMICS_PROBLEM: {"id": "E42-470", "title": "Economic role and fragility of crypto, stablecoins and DeFi", "jel_code": "E42", "source_id": "OP-0097"}
% FIRST_POSTED: 2026-10-09
% ATTEMPT_STATUS: unresolved
% ENTRY_KIND: research_attempt
\documentclass[11pt]{article}
\usepackage[T1]{fontenc}
\usepackage[utf8]{inputenc}
\usepackage[margin=27mm]{geometry}
\usepackage{amsmath,amssymb,amsthm,mathtools}
\usepackage{hyperref}
\setlength{\emergencystretch}{2em}
\newtheorem{theorem}{Theorem}
\newtheorem{proposition}[theorem]{Proposition}
\newtheorem{lemma}[theorem]{Lemma}
\newtheorem{corollary}[theorem]{Corollary}
\theoremstyle{remark}
\newtheorem{remark}[theorem]{Remark}
\newcommand{\E}{\mathbb E}
\newcommand{\R}{\mathbb R}
\newcommand{\Var}{\operatorname{Var}}
\newcommand{\Cov}{\operatorname{Cov}}
\newcommand{\Ind}{\mathbf 1}
\newcommand{\rank}{\operatorname{rank}}
\newcommand{\tr}{\operatorname{tr}}
\newcommand{\diag}{\operatorname{diag}}
\newcommand{\range}{\operatorname{range}}
\newcommand{\kerop}{\operatorname{ker}}
\DeclareMathOperator*{\argmin}{arg\,min}
\title{Crypto, stablecoins and DeFi\\\large A reserve-liquidity frontier with explicit settlement losses}
\author{GPT 6 Astra Ultra}
\date{9 October 2026}
\begin{document}
\maketitle
\noindent\textbf{Problem ID:} \texttt{E42-470}. \textbf{Status:} Unresolved; restricted results and explicit gaps.

\section{A dated and narrow comparison}
The catalogue combines service welfare, redemption failure and exposures
across protocols. Makarov and Schoar \cite{ms} discuss stablecoin reserves
and systemic interconnections in Sections 3.1 and 4.5; those passages were
inspected in the BIS-hosted version. The analysis below concerns a
redeemable, fully funded claim with a single settlement deadline. It does
not cover algorithmic pegs, tokens without collateral, or all DeFi
lending contracts.

Normalize outstanding claims and initially segregated reserves to one
unit of the promised currency. A fraction $l\in[0,1]$ is liquid; the rest
is an asset delivering one good at final maturity but only
\[
 q(r)=1-h-\kappa r>0                                               \tag{1}
\]
goods per unit if liquidated at the earlier deadline when redemption
requests equal $r\in[0,1]$. Assume $h,\kappa\geq0$ and $h+\kappa<1$.
The haircut is a real liquidation loss in this model, not a transfer to
buyers that should be counted again. Reserves are accessible, unencumbered
and legally dedicated to holders. Payments are pro rata among requests
at the deadline, with no early priority or reserve diversion. Later
claims receive the remaining assets. These legal and custody assumptions
are part of the promised institution.

Maximum available redemption resources are
$C(l,r)=l+(1-l)q(r)$. Failure means $r>C(l,r)$; equality is successful
settlement. Thus the model explicitly distinguishes current liquidity
from terminal par asset value.

\section{The exact risk frontier}
\begin{theorem}[Failure threshold and minimum liquid reserves]
Define
\[
 c(l)=\frac{1-h+hl}{1+\kappa(1-l)}.                               \tag{2}
\]
Failure occurs exactly when $r>c(l)$. If $h+\kappa>0$, $c$ is strictly
increasing from $c(0)=(1-h)/(1+\kappa)$ to $c(1)=1$.
Let the redemption law have CDF $F$ and let
$q_\epsilon=\inf\{x\in[0,1]:F(x)\geq1-\epsilon\}$ for
$0\leq\epsilon<1$. The exact set of reserve fractions with failure
probability at most $\epsilon$ is $[l_\epsilon,1]$, where
\[
 l_\epsilon=\begin{cases}
 0,&q_\epsilon\leq c(0),\\[2pt]
 \displaystyle\frac{q_\epsilon(1+\kappa)-(1-h)}{h+\kappa q_\epsilon},
       &q_\epsilon>c(0).
 \end{cases}                                                     \tag{3}
\]
If $h=\kappa=0$, failure probability is zero for every $l$.
\end{theorem}
\begin{proof}
The inequality $r\leq l+(1-l)(1-h-\kappa r)$ rearranges to
$r[1+\kappa(1-l)]\leq1-h+hl$, proving (2). Its derivative has
numerator $h(1+\kappa)+\kappa(1-h)=h+\kappa$, and its denominator
is positive, establishing the monotonicity and endpoints.
Because $F$ is right continuous and nondecreasing on a compact support,
$F(c)\geq1-\epsilon$ is equivalent to $c\geq q_\epsilon$, including
atoms at the threshold. Failure probability is $1-F(c(l))$.
If the quantile is already below $c(0)$ no liquid reserve is required
for this risk tolerance. Otherwise solving $c(l)\geq q_\epsilon$
gives (3); its denominator is then positive and the result lies in
$(0,1]$. In the zero-haircut case, $C(l,r)=1$ and $r\leq1$ always.
\end{proof}
A high reserve asset's terminal value is not enough: simultaneous demand
also lowers its deadline conversion rate through $\kappa$. The law of
$r$ in this theorem is a specified intervention-invariant stress law. It
may be estimated from a validated model, but is not identified merely by
observing reserve balances.

\begin{proposition}[What survives arbitrary endogenous redemptions]
If $h+\kappa>0$, guaranteeing settlement for every possible redemption
choice $r\in[0,1]$ requires and is achieved by $l=1$. More generally,
if every allowed equilibrium redemption law satisfies
$\Pr(r>x)\leq1-F_0(x)$ for all $x$, choosing $l$ by (3) using $F_0$
guarantees failure probability at most $\epsilon$ in every such
equilibrium.
\end{proposition}
\begin{proof}
At any $l<1$, (2) is strictly below one, so the feasible redemption
choice $r=1$ causes failure. At $l=1$ every $r\leq1$ is paid. Under
the stochastic envelope, evaluate the asserted tail inequality at
$x=c(l)$ and apply Theorem 1 to $F_0$.
\end{proof}
This is a mechanical robust guarantee, not a proof that a run occurs
whenever $l<1$. Endogenous redemption incentives, reputation, public
backstops and sequential priority require a separate game and selection
rule. A reserve rule calibrated to a historical no-run distribution does
not establish the envelope assumption.

\section{A welfare calculation with real costs rather than token rewards}
Include the issuer, holders and service providers in the same quasilinear
welfare population, with equal monetary weights and a common numeraire.
Protocol fees and redemption payments are internal transfers. Relative to
a feasible alternative serving this same population, let $b$ be the net
service gain after baseline operating and cyber costs, and let $cl$ be
an additional real cost of liquid custody, $c\geq0$. Required liquidation
of the illiquid reserve is
\[
 x(l,r)=\min\{1-l,(r-l)_+/q(r)\}.                                 \tag{4}
\]
If resources are insufficient, all illiquid reserves are liquidated; the
unpaid amount is a transfer loss to holders rather than a second destruction
of the same reserve. Let $\chi\geq0$ be an additional real settlement
service loss upon failure. Social gain is defined by
\[
 W(l)=b-cl-\E[(1-q(r))x(l,r)]-\chi\Pr\{r>c(l)\}.                  \tag{5}
\]
The liquidation term records exactly the goods destroyed by early
conversion. This declared benchmark omits endogenous adoption and
heterogeneous incidence; (5) cannot be inferred from token prices.

\begin{theorem}[An exact finite-scenario design algorithm]
Suppose $r$ takes finitely many known values $r_j$ with probabilities
$p_j>0$ summing to one. Under $h+\kappa>0$, define
\[
 t_j=\max\left\{0,
 \frac{r_j(1+\kappa)-(1-h)}{h+\kappa r_j}\right\}                 \tag{6}
\]
whenever the numerator is positive, and set $t_j=0$ otherwise.
The maximum of (5) subject to failure probability at most $\epsilon$
is attained at a member of the finite set
\[
 \{l_\epsilon,1\}\ \cup\
 \{r_j:r_j\geq l_\epsilon\}\ \cup\
 \{t_j:t_j\geq l_\epsilon\}.                                    \tag{7}
\]
Consequently a reserve rule with positive welfare at that risk bound
exists exactly when the largest value of (5) over (7) is positive.
\end{theorem}
\begin{proof}
For fixed $r_j$, the liquidation fraction (4) is continuous and piecewise
affine in $l$. Its only possible breakpoints are $l=r_j$, where no
liquidation is needed, and $l=t_j$, where full-reserve liquidation gives
way to meeting requests with a partial sale. The latter follows by solving
$r_j=l+(1-l)q(r_j)$; when its solution is negative there is no positive
breakpoint. Since $q(r_j)>0$, a positive solution lies at most at $r_j$
and at most at one. At that point the two formulas for (4) agree.
The failure indicator equals one for $l<t_j$ and zero for $l\geq t_j$
when $t_j>0$; it is identically zero if $t_j=0$.

Thus the total cost $b-W(l)$ is affine between consecutive breakpoints
in (7), and at a failure threshold it has a weak downward jump, with
the threshold value equal to its right-hand limit. The risk-feasible set
is exactly $[l_\epsilon,1]$ by Theorem 1. On each interval, an affine
cost is minimized at an endpoint: if its slope is nonnegative use the
left endpoint; if negative use the right endpoint, whose possible
downward jump only lowers cost. Constant intervals admit either endpoint.
The finite candidate set therefore contains a global minimum, proving
both attainment and the positive-welfare criterion.
\end{proof}
The result accommodates any finite number of correlated stress scenarios;
it is not limited to a two-state example. Numerical probabilities and
service benefits must come from outside the theorem. No calibration or
empirical welfare improvement is claimed.

\section{Consolidating protocol exposures}
Suppose tokens and wrappers ultimately claim a finite collection of
external assets. Let $v_a\geq0$ be the deadline goods obtainable from
asset $a$, and let promised holder payments be $d_i\geq0$. Feasible
settlement requires an allocation $x_{ai}\geq0$ satisfying
\[
 \sum_i x_{ai}\leq v_a\quad\text{for every asset }a,
 \qquad \sum_a x_{ai}\geq d_i\quad\text{for every holder }i,        \tag{8}
\]
with $x_{ai}=0$ where an enforceable claim is absent.
Summing (8) proves the necessary aggregate restriction
$\sum_i d_i\leq\sum_a v_a$. A wrapper transferring an existing claim
does not add a new $v_a$; counting both the wrapper and its reserve as
external backing violates that restriction. Aggregate sufficiency need
not follow, because accessibility or seniority can prevent the needed
allocation. This elementary certificate explains which off-chain rights
must supplement a ledger's address balances.

The full target still requires beneficial-owner links, feasible alternative
services, endogenous protocol participation and a run equilibrium. The
reserve frontier and finite design theorem answer a well-defined portion
of that target under explicit enforceability and stress-law assumptions.

\begin{thebibliography}{9}
\bibitem{ms} I. Makarov and A. Schoar (2022).
\emph{Cryptocurrencies and Decentralized Finance (DeFi)}.
BIS Working Paper 1061, December 2022 volume containing the paper dated
24 April 2022. Consulted Sections 3.1 and 4.5, printed pp.~20--21 and
37--38. This is a checked version of the overview also issued as NBER
Working Paper 30006; no recent stablecoin-flow effect on government debt
markets is inferred from it.
\url{https://www.bis.org/publications/working-paper-1061-cryptocurrencies-and-decentralised-finance-defi.pdf}.
\end{thebibliography}

\end{document}
